%%%%%%
%
% $Author: Abishek $
% $Date: 12.01.2025 $
% $Path: ML24-02-Gait-Pattern-Recognition\report\Contents\en$
% $Version: 1.0 $
%
%
% !TeX encoding = utf8
% !TeX root = Rename
% !TeX TXS-program:bibliography = txs:///bibtex
%
%%%%%%

\chapter{Algorithm}

\section{Description}
Convolutional Neural Networks (CNNs) are a class of deep learning algorithms primarily used for image and spatial data analysis. They are designed to automatically and adaptively learn spatial hierarchies of features from input images, making them highly effective for pattern recognition tasks like gait analysis. CNNs utilize layers such as convolutional layers, pooling layers, and fully connected layers to process and classify input data. \cite{TensorFlow:2023} 

\section{Applications}
CNNs are widely used in various domains, including:

\begin{itemize}
	\item \textbf{Image Classification:} Classifying images into categories, such as identifying objects or faces.
	\item \textbf{Object Detection:} Identifying the location of objects in an image.
	\item \textbf{Medical Image Analysis:} Recognizing patterns in medical imaging to detect conditions such as tumors or fractures.
	\item \textbf{Gait Recognition:} Classifying normal and abnormal walking patterns, as in the case of Parkinson’s disease, based on input data such as video or sensor data.
	\item \textbf{Speech and Audio Recognition:} Converting audio signals into meaningful patterns, such as speech-to-text.
\end{itemize}

\section{Relevance}
CNNs are highly relevant to gait pattern recognition because they excel at recognizing spatial and temporal patterns in data. In the context of gait analysis, CNNs can learn the complex features of human movement from sensor data or images, distinguishing between normal and Parkinsonian gaits. Their ability to process sequential data also helps in recognizing dynamic changes over time, which is essential in gait analysis.

\section{Hyperparameters}
The key hyperparameters for the CNN algorithm \cite{MNISTCNNgit:2023} include:

\begin{itemize}
	\item \textbf{Learning Rate:} Controls the step size during the optimization process. 
	\item \textbf{Batch Size:} Number of samples processed before the model’s internal parameters are updated.
	\item \textbf{Number of Epochs:} Number of times the learning algorithm will work through the entire training dataset.
	\item \textbf{Confidence Threshold:} Minimum confidence score for a detected object to be considered valid.
\end{itemize}

\section{Requirements}
The requirements for using CNNs in gait pattern recognition include:

\begin{itemize}
	\item \textbf{Hardware:} A powerful computing unit (e.g., Arduino Nano 33 BLE Sense Lite) for collecting sensor data.
	\item \textbf{Software:} Deep learning frameworks like TensorFlow or PyTorch for model training and deployment.
	\item \textbf{Dataset:} A labeled dataset with images or sensor readings representing various gait patterns (normal and Parkinsonian).
	\item \textbf{Preprocessing Tools:} Libraries for data cleaning, normalization, and feature extraction (e.g., Kalman filters for noise reduction).
\end{itemize}

\section{Input}
The input to the CNN algorithm consists of images captured by the Vision Shield. These images are preprocessed to grayscale and resized to a standard input size required by the model.

\section{Output}
The output of the CNN algorithm is the predicted class label for the input image, indicating which person it corresponds to based on the trained model. Additionally, the CNN provides confidence scores for each class, reflecting the model’s certainty about the prediction as a probability value between 0 and 1.

\section{Example with a program}

A simple example of a CNN for image classification, using the MNIST dataset (a dataset of handwritten digits) in Python with Tensorflow as shown in the code. \cite{Kaggle:2023}

\begin{lstlisting}[language=Python, caption=CNN using the MNIST dataset for image classification]
	
	import tensorflow as tf
	from tensorflow.keras import layers, models
	import numpy as np
	
	# Load the MNIST dataset (handwritten digits 0-9)
	(train_images, train_labels), (test_images, test_labels) = tf.keras.datasets.mnist.load_data()
	
	# Preprocess the data: Rescale pixel values to the range [0, 1] and reshape to (28, 28, 1)
	train_images = train_images / 255.0
	test_images = test_images / 255.0
	train_images = train_images.reshape(-1, 28, 28, 1)
	test_images = test_images.reshape(-1, 28, 28, 1)
	
	# Define the CNN model
	model = models.Sequential()
	
	# Add the first convolutional layer with 32 filters, kernel size (3, 3), and ReLU activation
	model.add(layers.Conv2D(32, (3, 3), activation='relu', input_shape=(28, 28, 1)))
	model.add(layers.MaxPooling2D((2, 2)))  # Max pooling layer to reduce spatial dimensions
	
	# Add the second convolutional layer with 64 filters
	model.add(layers.Conv2D(64, (3, 3), activation='relu'))
	model.add(layers.MaxPooling2D((2, 2)))
	
	# Flatten the output from convolutional layers
	model.add(layers.Flatten())
	
	# Add a fully connected layer with 64 units
	model.add(layers.Dense(64, activation='relu'))
	
	# Output layer with 10 units (one for each digit 0-9) and softmax activation for classification
	model.add(layers.Dense(10, activation='softmax'))
	
	# Compile the model with Adam optimizer and sparse categorical crossentropy loss function
	model.compile(optimizer='adam',
	loss='sparse_categorical_crossentropy',
	metrics=['accuracy'])
	
	# Train the model with training data
	model.fit(train_images, train_labels, epochs=5)
	
	# Evaluate the model on test data
	test_loss, test_accuracy = model.evaluate(test_images, test_labels)
	print(f'Test Accuracy: {test_accuracy * 100:.2f}%')
	
	# Make predictions on test data
	predictions = model.predict(test_images)
	print(f'Predictions for first test image: {np.argmax(predictions[0])}')
	
\end{lstlisting}